\section{Methods}
\label{sec:methods}

\subsection{What we are predicting}

Each row is one trial in one phase. A trial that spans two phases appears
twice, with the same success label, and those two rows always stay on the
same side of a split. Success means the trial met its primary endpoint.
The training file has 8,037 rows from trials that started through 2013. The
validation file has 1,013 rows from the same years, split by trial so a trial
is never in both. The test file has 1,500 rows, 500 in each phase, from 2013
to 2022. None of those trials appear in the other files.

Kaggle scores the average of the three phase log losses. Because each phase
has the same number of test rows, that is the same as the overall log loss.
We therefore select models by the equal average of the Phase~I, II, and III
log losses, not by a single F1 number. F1 at a cutoff of 0.5, ROC-AUC, and
precision-recall AUC are reported as diagnostics. Trial IDs are used only to
keep rows together and to match Part~1. They are not inputs. Part~1 reason
labels and reason embeddings are never given to the outcome model.

\subsection{The first models}

The baseline is a logistic regression on eligibility-criteria text. We turn
the criteria into word and bigram counts, keep at most 35,000 terms that
appear in at least three trials, and use sublinear term frequency. Three
phase indicators let one model share text evidence across phases. The
regularization is $C = 2$. A separate classifier for each phase was also
tried. Phase~I has many fewer rows, so a phase-only fit overfits it.

Around that baseline we compared disease and drug text, ICD codes, drug
targets, actions, and pathways~\cite{milacic2024reactome}, Morgan
fingerprints~\cite{rogers2010extended,rdkit}, naive Bayes, nearest neighbors,
and tree ensembles. Trees can represent interactions that a linear text model
misses. Fingerprints did not reliably help, and a matched drug is less common
in the test file (39\%) than in training (46\%), so a model that leans on drug
structure is exposed to a coverage shift. Naive Bayes could rank trials, but
its probabilities were far too sharp for a log-loss contest.

Our first submitted ensembles mixed logistic text models with CatBoost trees.
The version with the best local score put more weight on phase-specific
pieces and scored 0.59562 on Kaggle. Reducing the tree share in Phases~II
and~III improved that to 0.58993. The local winner was not the leaderboard
winner. That was the first sign that a score on the random validation split
was a weak guide to a later test cohort.

\subsection{The stack, and why recent trials count more}

The stronger reference is a stack of eight models~\cite{wolpert1992stacked,pedregosa2011scikit}.
Separate logistic regressions read eligibility text, disease text, drug names,
codes, drug annotations, and simple trial metadata. A naive Bayes model and a
nearest-neighbor model also read the eligibility text. A second logistic
regression learns how much to trust each of the eight. So that this combiner
does not see predictions that were fit on the answer, each base model is
scored on held-out folds: five folds, with every row of a trial kept together.
The neighbor model uses the 50 closest trials.

Trials get more successful over time, and every test trial starts later than
the training trials. We therefore give recent trials more weight. The weight
decays exponentially with a three-year time constant, so a trial loses about
half its weight every two years. The same weights go on the base models and
on the combiner. The decay was chosen on time-ordered holdouts, not on the
leaderboard. This recency-weighted stack, refit on all 9,050 labeled rows,
is the reported 0.58688 submission. An otherwise identical stack without the
recency weights scored 0.58810. Our first submissions had been fit on the
8,037 training rows only, so part of the earlier gap was the extra labeled
data, not only the architecture.

\subsection{Neighbors that notice time}

The neighbor piece inside that stack ignored recency. It used raw text
similarity and pulled every estimate toward a 50\% success rate. We tested
sharper similarity, neighborhoods of several sizes, similarity that also uses
disease, drug, and ICD text, and a prior estimated from the training rows of
the same phase.

Two versions mattered. One replaces only the neighbor input and then refits
the stack. The submitted neighbor version does not replace the original
stack. It averages that stack, at 50\%, with a second stack whose recency
decay is faster (a two-year time constant) and which also receives eight
neighborhood estimates. The average is fixed. It is not fit on the validation
labels. This is file~10. It improved every time-ordered window we had
declared in advance. Its random validation log loss got slightly worse.
Kaggle moved from 0.58688 to 0.58609.

\subsection{Why the submission mixes in the older ensemble}

File~10 was better on later years and worse on the random grouped checks.
Our older text-and-tree ensemble had the opposite pattern, so the two errors
were not the same. The older ensemble is half a mixture of three logistic
models, a quarter a depth-6 CatBoost model, and a quarter a depth-8 CatBoost
model. We blended probabilities rather than merging the features.

We tried 20\% of the older ensemble first, because a 30\% blend had a worse
worst-period score in the local checks. That 20\% blend, file~11, scored
0.58474. Raising the older share to 30\%, file~13, scored 0.58465. The second
gain was only 0.00009, so we did not keep adding more of the older model.
A later blend that also mixed in a recent tree and a phase trend looked
better on the historical windows and scored 0.58492, which is worse. With
one upload left and first place already reported, we did not submit the last
prepared file. Its local gains were small and its grouped score was slightly
worse.

The submitted probability is
\[
0.70\,P_{10} + 0.30\,P_{\text{early}},
\quad
P_{10} = 0.50\,P_{\text{recency}} + 0.50\,P_{\text{fast neighbors}}.
\]
After the recipe was fixed, every supervised piece was refit on all 9,050
labeled rows. Test outcome labels are never used. In the stack inherited from
the 0.58688 model, the word lists and IDF weights are fit on training,
validation, and unlabeled test text. That changes the features, not the
success labels. The older ensemble fits its own word lists only on the rows
it is allowed to learn from.

\subsection{How a candidate was allowed to advance}

Every later comparison uses the same saved trials. The supplied validation
split fits on 8,037 rows and scores 1,013. Three time-ordered windows train
through 2007, 2009, or 2011 and score the following years. A grouped
cross-validation pools the 9,050 labeled rows and holds out trials, twice,
with three folds. Three longer gaps were scored only after a recipe was
frozen. Word lists are rebuilt inside each fit, including the inner stacking
folds. We required a gain on the time-ordered average, and we recorded the
worst window and the grouped score even when they went the wrong way. A
higher F1 at 0.5 was not enough.

\subsection{The reason-error check}

After the outcome model is fit, we look up each failed validation trial in
Part~1. If Part~1 has a real reason label, we use it. Otherwise we use a
frozen Part~1 prediction. Trials that ran to completion are absent from
Part~1 and form their own ``no recorded stopping reason'' group. A trial with
several reasons is counted under each. Phases are separate, so a two-phase
trial is not counted twice inside one phase. The error rate is the share of
these failed trials given a success probability of at least 0.5. This table
is a diagnosis. It is not a feature and it is not the selection score.
